\section{Analiza zużycia pamięci RAM i obciążenia CPU}
\label{sec:RamCpu}

Kolejnym analizowanym aspektem było wykorzystanie zasobów sprzętowych podczas działania równoległych wariantów badanych algorytmów. 
Uwzględniono średnie obciążenie procesora oraz średnie i~maksymalne wykorzystanie pamięci RAM. 

W~celu bezpośredniego porównania wykorzystania zasobów przez badane algorytmy 
przyjęto wspólną konfigurację wykorzystującą 8 jednostek wykonawczych,
stanowiącą środkową wartość spośród badanych konfiguracji.

W tabeli~\ref{tab:resources} przedstawiono średnie obciążenie CPU oraz wykorzystanie pamięci RAM dla 
zbioru \texttt{random\_int} zawierającego 1~000~000 elementów.

\input{tables/chapter6/resources}

\subsection*{Obciążenie procesora CPU}

Wartości obciążenia CPU mogą przekraczać 100\%, ponieważ pomiar uwzględnia zarówno proces główny, 
jak i~procesy potomne wykorzystywane podczas wykonywania algorytmu równoległego.
W~zastosowanej metodzie pomiarowej 100\% oznacza pełne obciążenie jednego procesora logicznego, 
natomiast wyższe wartości wskazują na równoczesne wykorzystanie kilku jednostek wykonawczych.

Przy wykorzystaniu 8 jednostek wykonawczych najwyższe średnie obciążenie CPU odnotowano dla Sample Sort i~wyniosło ono około 426,09\%.
Drugą najwyższą wartość uzyskał Bucket Sort z wynikiem około 365,52\%. 
Merge Sort oraz Quick Sort charakteryzowały się wyraźnie mniejszym obciążeniem procesora, osiągając odpowiednio około 180,70\% i~173,84\%.

Zmiany obciążenia procesora wraz ze wzrostem liczby jednostek wykonawczych przedstawiono na rysunku~\ref{fig:cpu-cores-random-int}.
Dodatkowo zaznaczono teoretyczne maksymalne obciążenie CPU, odpowiadające 100\% dla każdej wykorzystywanej jednostki wykonawczej.

\begin{figure}[htbp]
    \centering
    \includegraphics[width=0.95\textwidth]
    {images/chapter6/cpu/cpu_vs_cores_comparison_random_int.png}
    \caption{Średnie wykorzystanie CPU dla zbioru \texttt{random\_int}
    o rozmiarze 1~000~000 elementów w~zależności od liczby jednostek
    wykonawczych. Źródło: opracowanie własne.}
    \label{fig:cpu-cores-random-int}
\end{figure}

Dla Bucket Sort zwiększanie liczby jednostek wykonawczych skutkowało stopniowym wzrostem obciążenia procesora. 
Średnie wykorzystanie CPU wzrosło z około 163,59\% dla 2 jednostek do 232,20\% dla 4,
365,52\% dla 8 oraz 529,63\% dla 16 jednostek wykonawczych.
Największe średnie obciążenie CPU odnotowano przy 32 jednostkach wykonawczych i~wyniosło ono około 633,31\%.
Wzrost obciążenia CPU nie prowadził jednak do dalszego skracania czasu sortowania dla wszystkich badanych konfiguracji.
Najlepszy czas dla tego algorytmu uzyskano przy 16 jednostkach wykonawczych, 
natomiast zwiększenie ich liczby do 32 spowodowało wydłużenie czasu do około 1,48~s, pomimo dalszego wzrostu średniego wykorzystania procesora. 
Wynik ten pokazuje, że samo zwiększenie obciążenia CPU nie prowadziło do poprawy czasu wykonania.

Odmienną tendencję zaobserwowano dla Merge Sort. 
Wykorzystanie CPU wzrosło z około 181,93\% dla 2 jednostek wykonawczych do około 253,66\% dla 4 jednostek, 
po czym przy 8 jednostkach spadło do około 180,70\%. 
Dla 16 oraz 32 jednostek wykonawczych średnie obciążenie procesora utrzymywało się na zbliżonym poziomie, 
wynosząc odpowiednio około 183,61\% i~182,38\%.
Zależność ta odpowiada wynikom analizy czasu sortowania przedstawionym w~sekcji~\ref{sec:LiczbaRdzeni}, 
gdzie zwiększanie liczby jednostek wykonawczych powyżej 4 nie prowadziło do istotnego skrócenia czasu wykonania.

Podobną zależność zaobserwowano dla Quick Sort. 
Dla 2 jednostek wykonawczych średnie obciążenie procesora wynosiło około 173,30\%, natomiast dla 4 jednostek wzrosło do około 215,68\%. 
Po zwiększeniu liczby jednostek do 8 obciążenie spadło do około 173,84\%, a~dla 16 i~32 jednostek wynosiło odpowiednio około 171,53\% i~175,47\%.
Pomimo zwiększania liczby jednostek wykonawczych nie obserwowano zatem systematycznego wzrostu wykorzystania procesora. 
Jest to zgodne z~wynikami czasu wykonania, dla których zwiększanie poziomu równoległości powyżej 2 jednostek nie przynosiło poprawy.

Wyraźne zmiany wystąpiły również dla Sample Sort. 
Średnie obciążenie procesora wzrosło z około 190,01\% dla 2 jednostek wykonawczych do 297,86\% dla 4 oraz 426,09\% dla 8 jednostek.
Dalsze zwiększenie liczby jednostek powodowało dalszy wzrost obciążenia procesora, 
które wyniosło około 565,48\% dla 16 oraz 609,04\% dla 32 jednostek wykonawczych.

Porównanie wszystkich algorytmów wskazuje zatem na wyraźne różnice w~sposobie wykorzystania dostępnej równoległości. 
Bucket Sort oraz Sample Sort wykazywały wzrost średniego obciążenia CPU wraz ze zwiększaniem liczby jednostek wykonawczych. 
W~przypadku Bucket Sort wzrost ten był szczególnie wyraźny w~całym badanym zakresie, 
natomiast Sample Sort osiągał wysokie wartości obciążenia zwłaszcza dla 16 i~32 jednostek. 
Merge Sort i~Quick Sort charakteryzowały się natomiast znacznie niższym oraz bardziej stabilnym wykorzystaniem procesora 
po przekroczeniu 4 jednostek wykonawczych.

Kolejnym czynnikiem wpływającym na wykorzystanie procesora był rozmiar przetwarzanego zbioru. 
Zależność średniego obciążenia CPU od liczby sortowanych elementów dla 32 jednostek wykonawczych 
przedstawiono na rysunku~\ref{fig:cpu-size-random-int}.

\begin{figure}[htbp]
    \centering
    \includegraphics[width=0.95\textwidth]
    {images/chapter6/cpu/cpu_vs_data_size_random_int.png}
    \caption{Średnie wykorzystanie CPU w~zależności od rozmiaru zbioru
    \texttt{random\_int} dla 32 jednostek wykonawczych. Źródło: opracowanie własne.}
    \label{fig:cpu-size-random-int}
\end{figure}

Dla najmniejszych zbiorów danych w części pomiarów odnotowano bardzo niskie,
a~w~niektórych przypadkach również zerowe wartości średniego wykorzystania CPU.
Wyniki te należy interpretować z~uwzględnieniem przyjętego sposobu próbkowania obciążenia procesora.
Ze względu na krótki czas wykonania sortowania dla najmniejszych zbiorów liczba zarejestrowanych próbek była niewielka. 
Ograniczało to dokładność oceny średniego poziomu obciążenia CPU dla tych konfiguracji.
Wartość równa 0\% nie oznacza zatem braku wykorzystania procesora, lecz wynika z~ograniczenia zastosowanej metody pomiarowej.
Dla większych zbiorów danych rejestrowane wartości wykorzystania procesora były wyższe. 
Dla zbioru \texttt{random\_int} o rozmiarze 100~000 elementów średnie wykorzystanie procesora przy 
32 jednostkach wykonawczych wynosiło od około 65\% do 86\%, zależnie od algorytmu. 
Po~zwiększeniu rozmiaru zbioru do 1~000~000 elementów wartości te wzrosły wyraźnie, 
osiągając około 175--182\% dla Quick Sort i Merge Sort oraz ponad 600\% dla Bucket Sort i Sample Sort.

Szczególnie wyraźny wzrost obciążenia procesora zaobserwowano dla Bucket Sort oraz Sample Sort. 
Dla obu algorytmów zwiększenie rozmiaru danych z~100~000 do 1~000~000 elementów prowadziło do wielokrotnego wzrostu średniego wykorzystania CPU. 
W~przypadku Merge Sort i~Quick Sort wzrost był znacznie mniejszy. 
Wskazuje to, że nie tylko liczba jednostek wykonawczych, 
ale również rozmiar przetwarzanego zbioru wpływał na możliwość efektywnego wykorzystania dostępnych zasobów.

Zbliżone zależności zaobserwowano również dla zbioru \texttt{random\_float}.
Dla 32 jednostek wykonawczych średnie obciążenie CPU dla 1~000~000 elementów wyniosło 
około 657,02\% dla Bucket Sort, 182,29\% dla Merge Sort, 136,53\% dla Quick Sort oraz 650,69\% dla Sample Sort. 
Również w~tym przypadku Bucket Sort i~Sample Sort wykazywały znacznie większe wykorzystanie procesora niż Merge Sort i~Quick Sort. 
Potwierdza to, że zaobserwowane zależności nie były charakterystyczne wyłącznie dla danych typu całkowitego. 

\subsection*{Wykorzystanie pamięci RAM}

Dla przyjętej wcześniej wspólnej konfiguracji 8 jednostek wykonawczych przeanalizowano również wykorzystanie pamięci RAM przez poszczególne algorytmy.
W~przypadku zbioru \texttt{random\_int} największe średnie zużycie pamięci odnotowano dla Sample Sort i~wyniosło ono około 258,15~MB.
Zbliżony wynik uzyskał Quick Sort z wartością około 249,93~MB, natomiast dla Bucket Sort było to około 223,00~MB. 
Najniższe średnie wykorzystanie pamięci zarejestrowano dla Merge Sort i~wyniosło około 163,15~MB.

Pod względem maksymalnego zużycia pamięci najwyższą wartość odnotowano dla Sample Sort i~wyniosła ona około 480,09~MB.
Dla Bucket Sort maksymalne wykorzystanie pamięci wyniosło około 442,46~MB, dla Quick Sort około 366,78~MB, natomiast dla Merge Sort około 216,79~MB.
Oznacza to, że przy tej samej liczbie jednostek wykonawczych poszczególne 
algorytmy charakteryzowały się wyraźnie różnym zapotrzebowaniem na pamięć operacyjną.

Wpływ zwiększania liczby jednostek wykonawczych na wykorzystanie pamięci RAM przedstawiono na rysunku~\ref{fig:memory-cores-random-int}.

\begin{figure}[htbp]
    \centering
    \includegraphics[width=0.95\textwidth]
    {images/chapter6/memory/memory_vs_cores_comparison_random_int.png}
    \caption{Średnie i maksymalne wykorzystanie pamięci RAM dla zbioru
    \texttt{random\_int} o rozmiarze 1~000~000 elementów w~zależności od liczby jednostek wykonawczych. Źródło: opracowanie własne.}
    \label{fig:memory-cores-random-int}
\end{figure}

Najbardziej wyraźny wzrost wykorzystania pamięci wraz z~liczbą jednostek wykonawczych wystąpił dla Bucket Sort. 
Dla zbioru zawierającego 1~000~000 elementów średnie wykorzystanie pamięci wzrosło z~około 175,02~MB 
dla 2 jednostek wykonawczych do 187,12~MB dla 4, 223,00~MB dla 8, 275,18~MB dla 16 oraz 327,15~MB dla 32 jednostek wykonawczych.
Jednocześnie maksymalne wykorzystanie pamięci wzrosło od około 230,50~MB dla 2 jednostek do około 1~182,75~MB dla 32 jednostek.

Wzrost wykorzystania pamięci dla Bucket Sort był związany ze zwiększaniem poziomu równoległości. 
Jednocześnie zwiększenie liczby jednostek wykonawczych z~16 do 32 nie przyniosło skrócenia czasu sortowania, 
lecz doprowadziło do dalszego wzrostu zużycia pamięci.
Zwiększenie poziomu równoległości wiązało się z~dodatkowym kosztem pamięciowym,
któremu w~przypadku zwiększenia liczby jednostek z~16 do 32 towarzyszyło wydłużenie czasu wykonania.

W przypadku Merge Sort zależność między liczbą jednostek wykonawczych a zużyciem pamięci była mniej wyraźna. 
Najwyższe średnie wykorzystanie pamięci odnotowano przy 4 jednostkach wykonawczych, natomiast przy 8 jednostkach nastąpił spadek do poziomu około 163~MB. 
Dalsze zwiększanie liczby jednostek do 16 i 32 nie powodowało już istotnych zmian, a wartości pozostawały na zbliżonym poziomie. 
Maksymalne wykorzystanie pamięci również nie wykazywało jednoznacznej tendencji wzrostowej wraz ze zwiększaniem liczby jednostek wykonawczych.

W przypadku Quick Sort średnie wykorzystanie pamięci wzrosło z około 200~MB dla 2 jednostek do około 263~MB dla 4 jednostek wykonawczych. 
Przy dalszym zwiększaniu ich liczby wartości pozostawały już na zbliżonym poziomie, wynosząc około 250--253~MB. 
Maksymalne wykorzystanie pamięci było najwyższe przy 4 jednostkach i wyniosło około 433~MB, 
natomiast dla większej liczby jednostek spadło do około 366--371~MB. Oznacza to, że podobnie jak w przypadku średniego wykorzystania pamięci, 
zależność między liczbą jednostek wykonawczych a zużyciem pamięci nie była monotoniczna.

Dla Sample Sort zaobserwowano wyraźny wzrost średniego wykorzystania pamięci wraz ze zwiększaniem liczby jednostek wykonawczych. 
Wartość ta zwiększała się z około 204~MB dla 2 jednostek do około 326~MB dla 32 jednostek. 
Jeszcze wyraźniejszy wzrost wystąpił w przypadku maksymalnego wykorzystania pamięci, które zwiększyło się z około 272~MB do ponad 1,2~GB. 
Oznacza to, że w przypadku Sample Sort zwiększanie liczby jednostek wykonawczych wiązało się z wyraźnym wzrostem zapotrzebowania na pamięć.

Porównanie obu typów danych wskazuje na zbliżone wykorzystanie pamięci RAM. 
Dla obu zbiorów średnie zużycie pamięci przy 8 jednostkach wykonawczych było na podobnym poziomie, 
co wskazuje, że typ danych nie wpływał istotnie na charakterystykę wykorzystania pamięci w~badanej konfiguracji.

Większe różnice zaobserwowano w przypadku maksymalnego wykorzystania pamięci przy 32 jednostkach wykonawczych. 
Dla obu typów danych najwyższe wartości wystąpiły dla Bucket Sort oraz Sample Sort i przekraczały 1~GB. 
Różnice pomiędzy \texttt{random\_int} i \texttt{random\_float} były jednak niewielkie, co potwierdza, 
że w~badanej konfiguracji typ danych nie zmieniał zasadniczo charakterystyki wykorzystania pamięci przez poszczególne algorytmy.